\section*{Chapter \ref{ch:s1:earnings} --- Earnings}

\begin{solution}{exo:s1:earnings:1}
$1.5 \times 1.2\% = 1.8\%$ over sixty days; per unit of surprise $1.2\%/60 = 2.0$ basis points a day.
\end{solution}

\begin{solution}{exo:s1:earnings:2}
The weight traded is twice the one-way turnover: $2 \times 65.2 \times 0.001 = 13.0\%$ a year for the five-day book, $2 \times 6.6 \times 0.001 = 1.3\%$ for
the sixty-day book. The five-day book's gross return, 8.3\% a year, does not cover its costs.
\end{solution}

\begin{solution}{exo:s1:earnings:3}
About $0.05 \times 6.94\% = 0.35$ points, close to the 0.33 the \omterm{def:s1:earnings:event}{event study} shows (the move also contains market and industry parts, which are not reversed).
\end{solution}

\begin{solution}{exo:s1:earnings:4}
The announcement-day return is the surprise's jump plus that day's specific noise, so it measures the surprise with error; the reported surprise is the
planted cause itself. On real data the reported surprise is measured against a consensus that may be stale or biased, and the return reflects everything
announced (guidance, margins), so the return can carry information the headline number lacks. Combining both is common.
\end{solution}

\begin{solution}{exo:s1:earnings:5}
The date is scheduled and published in advance; the surprise becomes public only at the announcement. A backtest must know whether each announcement came
before the open, during the session or after the close, and act only from the first price after it: an after-close announcement cannot be traded at that
day's close.
\end{solution}

\begin{solution}{exo:s1:earnings:6}
Over five trading days with one announcement day whose variance is about nine times an ordinary day's ($3.02^2$), the variance is $4 + 9 = 13$ ordinary days'
against 5 without the event: the implied volatility of the option spanning the announcement is about $\sqrt{13/5} = 1.61$ times that of one that does
not, and falls back once the event is past.
\end{solution}

\begin{solution}{exo:s1:earnings:7}
Sharpe ratios of 2.18 before and 0.75 after the break; 4.8\% a year before and 1.4\% after. A drop of that size in a series with a volatility of about 2\% a
year takes a few years to separate from noise with Book~7, chapter~13's break tests; waiting for certainty means trading a weakened strategy for years.
\end{solution}

\begin{solution}{exo:s1:earnings:8}
The drift is concentrated in illiquid stocks: Chordia and co-authors found 0.04\% a month in the most liquid stocks against 2.43\% in the most illiquid, and
Martineau none in large stocks since 2006. A backtest on all stocks says little about the S\&P 500; run it on that universe, with costs, and on recent
years.
\end{solution}

\begin{solution}{pb:s1:earnings:1}
\begin{enumerate}
\item Average abnormal returns in \omterm{def:s1:earnings:event}{event time} over a window around events; the drift is the continuation of abnormal returns in the surprise's direction.
\item Past surprises and past returns each predict large drifts controlling for the other; analysts' forecasts also respond sluggishly.
\item A standard normal surprise, a jump of three daily specific volatilities per unit and 1.2\% of drift per unit over sixty days; four announcements a year
(1.59\% of stock-days).
\item Above $+1$ s.d.: 6.94\% on the day, 7.75\% at sixty days; below $-1$: $-6.93\%$ and $-8.56\%$.
\item Long names with a surprise above one standard deviation and short those below minus one, from a delay after the announcement for a holding window,
gross one.
\item At the announcement's close: $-3.72$, 0.02, 2.03; a day later: $-0.77$, 1.51, 2.66.
\item The next day reverses part of the jump (0.33 points for the positive group).
\item Each position pays a full round trip for a few basis points of drift a day: costs of 13\% a year at five days.
\item The surprise (2.66 against 2.03 at sixty days): the return mixes the surprise with noise.
\item Higher returns around scheduled announcements, 9.9\% a year annualised in Savor and Wilson; none on the synthetic market (1.8 basis points a day, $t =
0.59$).
\item 3.02 times ordinary moves; options spanning the date price the jump, so their implied volatility is higher.
\item Buying upward revisions of forecasts, and trading on guidance and pre-announcements; point-in-time forecast histories and timestamped press releases.
\item \textbf{Named result.} Sharpe ratios after costs of $-0.77$, 1.51 and 2.66 at holding periods of 5, 20 and 60 days (bought a day after the
announcement); after the drift loses 70\% of its size, the sixty-day book's Sharpe ratio over years 7--10 is 0.75 against 3.24.
\item It is weak or absent in liquid stocks (0.04\% a month) and absent in large stocks since 2006.
\item In small, illiquid stocks, where costs eat it, and in signals the market digests more slowly (guidance, revisions, text).
\item Track realised drift per unit of surprise, its event-study path, and the book's return against a break test.
\item Before-or-after-market timestamps, stale consensus, restated earnings, and announcement dates moved after the fact.
\item The sixty-day drift book in the liquid names where it survives, small; or the \omterm{def:s1:earnings:premium}{announcement premium} as a risk premium.
\item That a return concentrated around events can be a reward for risk, not only a mistake.
\item Positions taken around scheduled news, betting on how completely prices absorb it.
\end{enumerate}
\end{solution}

\begin{solution}{iq:s1:earnings:1}
Abnormal returns continuing in the direction of an earnings surprise for weeks after the announcement. Explanations: investors underreact and update slowly;
limits to arbitrage in the illiquid stocks where it is strongest; a risk premium the literature has not settled.
\end{solution}

\begin{solution}{iq:s1:earnings:2}
Actual earnings as first reported and the consensus as it stood before the announcement, both timestamped; the surprise scaled by price or by the history
of surprises (SUE). What goes wrong: restated earnings replacing first reports, consensus snapshots taken after the announcement, announcement times
misdated, and split-adjusted per-share numbers mixed with unadjusted ones.
\end{solution}

\begin{solution}{iq:s1:earnings:3}
Test for a break formally, estimate its size after the break, measure it where the book can trade (liquid names, with costs), and ask why it weakened. If
the post-break return still exceeds costs with margin, keep a smaller allocation; if not, stop.
\end{solution}

\begin{solution}{iq:s1:earnings:4}
Compare the move implied by the options spanning the date (from the difference between their implied variance and the surrounding options') with a
forecast of the actual move; buy delta-hedged straddles when the forecast exceeds the implied move by more than costs, sell them in the opposite case.
\end{solution}

\begin{solution}{iq:s1:earnings:5}
Jump risk on every announcement, concentration in the names announcing that week, earnings-season clustering, and data risk (timestamps, restatements).
\end{solution}

\begin{solution}{iq:s1:earnings:6}
With positions held $H$ days and replaced, there are $252/H$ round trips a year per unit of capital. For $H \le 60$ the net return is
$(252/H)(dH/60 - 2c) = 252\,(d/60 - 2c/H)$, increasing in $H$; for $H > 60$ it is $(252/H)(d - 2c)$, decreasing. The optimum is to hold for the whole drift,
$H = 60$.
\end{solution}
